%%% =========================================================
%%% Sprache, Encoding, Fonts
%%% =========================================================
\usepackage[utf8]{inputenc}
\usepackage[T1]{fontenc}
\usepackage[ngerman]{babel}
\usepackage{lmodern}

%%% =========================================================
%%% Mathematik & Symbole
%%% =========================================================
\usepackage{amsmath}
\usepackage{amsfonts}
\usepackage{amssymb}
\usepackage{dsfont}

%%% =========================================================
%%% Tabellen, Layout, Farben
%%% =========================================================
\usepackage{array}
\usepackage{tabularx}
\usepackage{longtable}
\usepackage{ltablex} % Tipp: \keepXColumns nach \begin{document}
\usepackage{booktabs}
\usepackage{dcolumn}
\usepackage[table,xcdraw]{xcolor}
\usepackage{nameref}

%%% =========================================================
%%% Seitenlayout & Ränder
%%% =========================================================
\usepackage[top=2.5cm, bottom=2.5cm, left=2cm, right=2cm, head=17pt]{geometry}
\usepackage{pdflscape}

%%% =========================================================
%%% Typografie, Abstände, Absätze
%%% =========================================================
\usepackage{microtype}
\usepackage{setspace} % Verwende \onehalfspacing im Dokumentteil
\usepackage{blindtext}
\usepackage{lipsum}
\usepackage{titlesec}
\usepackage{enumitem}
\usepackage{lettrine}

%%% Absatzformatierung
\titleformat{\paragraph}{\normalfont\normalsize\bfseries}{\theparagraph}{1em}{}
\titlespacing*{\paragraph}{0pt}{3.25ex plus 1ex minus .2ex}{1.5ex plus .2ex}

%%% =========================================================
%%% Grafiken & Floats
%%% =========================================================
\usepackage{graphicx}
\usepackage{subcaption}
\usepackage{adjustbox}
% \usepackage{floatflt} % obsolet -> ggf. durch wrapfig/floatrow ersetzen
\usepackage{epstopdf}
\usepackage{float} % statt 'here'; ermöglicht [H]

%%% =========================================================
%%% Listings / Code
%%% =========================================================
\usepackage{listings}
\usepackage{etoolbox} % für \ifstrempty

%%% =========================================================
%%% Datum & Zeit
%%% =========================================================
\usepackage[ngerman]{datetime2}
\newcommand{\timestamp}{\DTMcurrenttime}
\DTMsetdatestyle{ddmmyyyy}
%\DTMsetup{datesep={.}}

\usepackage{advdate}


%%% =========================================================
%%% Kopf- und Fußzeilen
%%% =========================================================
\usepackage[
  headsepline=0.4pt,
  footsepline=0.4pt,
  plainheadsepline,
  plainfootsepline
]{scrlayer-scrpage}

\usepackage[
  stable,    % Fußnoten in Überschriften erlauben
  multiple,  % Komma zwischen Fußnoten (1,2)
  perpage,   % Zähler auf jeder Seite bei 1 beginnen
  bottom     % Fußnoten immer ganz unten fixieren
]{footmisc}

\usepackage{tablefootnote} % Für Fußnoten in Tabellen

% Fußnote in kleinem römischen Zahlen
\renewcommand{\thefootnote}{\roman{footnote}}

%%% =========================================================
%%% Glossare & Akronyme
%%% =========================================================
\usepackage[acronym, toc, nomain, section=subsection]{glossaries}
\usepackage{xparse}
\makenoidxglossaries

\newglossary*{general}{Glossar}

\NewDocumentCommand{\myacronym}{m m m O{}}{%
  % Akronym anlegen; Beschreibung in user1 speichern
  \newacronym[user1={#4}]{#1}{#2}{#3}%
  % Falls Beschreibung vorhanden → zusätzlichen Glossar-Eintrag im Typ "general"
  \ifstrempty{#4}{}{%
    \newglossaryentry{gls-#1}{%
      type=general,
      name={#2},
      description={#3\textemdash\ #4},
      user1={#4}% wichtig: damit dein Glossar-Style ihn nicht wegfiltert
    }%
  }%
}

%%% ---------------------------------------------------------
%%% Glossar-Stil: zweispaltig_akk (Abkürzungsverzeichnis)
%%% ---------------------------------------------------------
\newglossarystyle{zweispaltig_akk}{%
  \renewenvironment{theglossary}%
  {%
    \renewcommand{\arraystretch}{1.2}%
    \begin{longtable}{|
      p{3cm}|
      p{\dimexpr\textwidth-3cm-4\tabcolsep-3\arrayrulewidth\relax}|
    }
    \hline
  }{%
    \caption{Abkürzungsverzeichnis}%
    \label{tab:akronyme}%
    \\ \end{longtable}
  }%

  \renewcommand*{\glossaryheader}{%
    \rowcolor[HTML]{9B9B9B}
    \textbf{Akronym} & \textbf{Langform} \\ \hline
    \endhead
  }%

  \renewcommand*{\glossentry}[2]{%
    \glstarget{##1}{\glossentryname{##1}} &
    \glossentrydesc{##1} \\
    \hline
  }%

  \renewcommand*{\glsgroupskip}{}%
}

%%% ---------------------------------------------------------
%%% Glossar-Stil: glossar_zweispaltig (Glossar)
%%% ---------------------------------------------------------
\newglossarystyle{glossar_zweispaltig}{%
  \renewenvironment{theglossary}%
  {\renewcommand{\arraystretch}{1.2}%
   \begin{longtable}{|p{3cm}|p{\dimexpr\textwidth-3cm-4\tabcolsep-3\arrayrulewidth\relax}|} \hline}%
  {%
    \caption{Glossar}
    \label{tab:mein_glossar}
    \\ \end{longtable}}%
  \renewcommand*{\glossaryheader}{%
    \rowcolor[HTML]{9B9B9B}
    \textbf{Begriff} & \textbf{Erläuterung} \\ \hline \hline \endhead}%
  % Nur drucken, wenn user1 (=\glsentryuseri) NICHT leer ist:
  \renewcommand*{\glossentry}[2]{%
    \ifstrempty{\glsentryuseri{##1}}%
      {}%
      {%
        \glstarget{##1}{\glossentryname{##1}} &
        {\textit{\glsentryuseri{##1}}} \\ \hline
      }%
  }%
  \renewcommand*{\glsgroupskip}{}%
}

%%% =========================================================
%%% Wiederholbare Abbildungs-Liste
%%% =========================================================
\usepackage{atveryend}

\makeatletter
\newcommand\Repeatable@starttoc[1]{%
  \begingroup
  \makeatletter
  \@input{\jobname.#1}%
  \if@filesw
    \AfterLastShipout{%
      \@ifundefined{tf@#1}{%
        \expandafter\newwrite\csname tf@#1\endcsname
        \immediate\openout \csname tf@#1\endcsname \jobname.#1\relax
      }{}%
    }%
  \fi
  \@nobreakfalse
  \endgroup
}%
\newcommand\listoffiguresRepeatable{%
  \begingroup
  \let\@starttoc\Repeatable@starttoc
  \listoffigures
  \endgroup
}%
\makeatother


%%% =========================================================
%%% Datumsrechnerei
%%% =========================================================


%%% =========================================================
%%% Fußnoten / Endnoten / Kommentare
%%% =========================================================
\usepackage{comment}

%%% Endnoten referenzierbar machen
\usepackage{etoolbox}
\usepackage{enotez}

\setenotez{
  list-name      = {Endnoten},
  totoc          = false,
  counter-format = {arabic},
  backref        = true,
  reset          = true,
  list-style     = plain
}

%%% =========================================================
%%% PDF-Seiten einbinden
%%% =========================================================
\usepackage{pdfpages}

%%% =========================================================
%%% Hyperref (immer vor cleveref)
%%% =========================================================
\usepackage[pdfborderstyle={/S/U/W 1}]{hyperref}

%%% =========================================================
%%% Cleveref (immer nach hyperref!)
%%% =========================================================
\usepackage{cleveref}

%%% =========================================================
%%% PDF-Metadaten
%%% =========================================================
\hypersetup{
    pdftitle={Report Sebastian Seitz - IPMA Level B},
    pdfauthor={Sebastian Seitz},
    pdfsubject={Governance, Strukturen und Prozesse},
    % --- Link-Styling ---
    colorlinks=true,    % Links einfärben statt Rahmen
    linkcolor=black,    % Interne Links (Inhaltsverzeichnis, Endnoten)
    citecolor=black,    % Zitate / Literatur
    filecolor=black,    % Dateilinks
    urlcolor=blue,      % Web-Adressen (hier ist Blau sinnvoll)
    bookmarksnumbered=true
}


\usepackage{fmtcount}